% Einheit 4 von 4 – Integralfunktion (bank/stammfunktion-und-hauptsatz/e4.jsonl, zusammenbau v0.1)
%% TODO Zeitmarke Einheit 4: Marken-Zeile nicht lesbar
%% TODO Zweigzeile Teil 1: was hier gelernt wird (Satz in Schülersprache aus dem Einheitstitel) – Einheit 4
\einheitenkopf[e4]{Einheit 4 von 4 · Integralfunktion}
\zweigzeile{keine Prüfungsaufgabe}
\setcounter{aufgabe}{14}

%% TODO Titel: Ich-kann-Satz fehlt in der Bank (Platzhalter: Kettenname) – Nr. 15
%% TODO Anweisung: ein Satz über der ersten Teilaufgabe fehlt in der Bank – Nr. 15
\begin{aufgabe}{Integralfunktion}
%% TODO \rechenplatz{n}: Schrittzahl des Grundfalls fehlt in der Bank (nur bei fester Schreibform, 2.3 b)
\begin{teile}
\teil Das Bild zeigt den Graphen von $f$. Ist das Integral von $-2$ bis $2$ über $f(x)$ null? \\ \kreuz{ja – gleiche Grenzen} \\ \kreuz{ja – die Flächenstücke heben sich auf} \\ \kreuz{nein}

\begin{ksys}[xmin=-3,xmax=3,ymin=-3,ymax=3,ablesen] \gerade{1}{0}{f} \end{ksys}
\teil $J(x)$ ist das Integral von $2$ bis $x$ über $f(t) = t^2 - 1$. Gib eine Nullstelle von $J$ und den Term von $J'$ an. Nullstelle: \leerfeld , $J'(x) = $ \leerfeld
\teil Das Bild zeigt den Graphen von $f$; er ist punktsymmetrisch zu $(3 | 0)$. $J(x)$ ist das Integral von $1$ bis $x$ über $f(t)$. Gib die beiden ganzzahligen Nullstellen von $J$ mit $0 \le x \le 6$ an und begründe. \hfill (Abitur 2018 LK)

\begin{ksys}[xmin=-1,xmax=7,ymin=-4,ymax=4,ablesen] \funktionab{0.25*\x*(\x-3)*(\x-6)}{f}{-0.5}{6.5} \end{ksys}
\teil $J(x)$ ist das Integral von $0$ bis $x$ über $f(t) = t^2 - 4$. Wie viele Nullstellen kann $J$ höchstens haben? Begründe. höchstens \leerfeld
\teil $J(x)$ ist das Integral von $0$ bis $x$ über $f(t) = t^2 - 4t$. Bestimme die Wendestelle von $J$ und den Funktionswert dort. Wendestelle: \leerfeld , Wert: \leerfeld
\teil Das Bild zeigt den Graphen von $f$ mit den Nullstellen $1$ und $3$; für $x > 3$ fällt er unbeschränkt. $J(x)$ ist das Integral von $1$ bis $x$ über $f(t)$. Aussage: $J$ hat genau zwei Nullstellen. Beurteile die Aussage. \hfill (Abitur 2026 LK) \\ \kreuz{wahr} \\ \kreuz{falsch}

\begin{ksys}[xmin=-1,xmax=5,ymin=-6,ymax=4,ablesen] \funktionab{-0.5*(\x-1)^2*(\x-3)}{f}{-0.5}{4.1} \end{ksys}
\end{teile}
\end{aufgabe}

%% TODO Titel: Ich-kann-Satz fehlt in der Bank (Platzhalter: Kettenname) – Nr. 16
%% TODO Anweisung: ein Satz über der ersten Teilaufgabe fehlt in der Bank – Nr. 16
\begin{aufgabe}{Integralfunktionen mit verschiedenen unteren Grenzen als Verschiebung in y-Richtung begründen}
\begin{teile}
\teil $f(t) = 4t - t^2$. $L(x)$ ist das Integral von $0$ bis $x$ über $f(t)$, $M(x)$ das Integral von $3$ bis $x$ über $f(t)$. Begründe, dass der Graph von $M$ aus dem von $L$ durch eine Verschiebung in negative y-Richtung hervorgeht, und gib an, um wie viel. \hfill (Abitur 2017 LK)
\end{teile}
\end{aufgabe}

%% TODO Titel: Ich-kann-Satz fehlt in der Bank (Platzhalter: Kettenname) – Nr. 17
%% TODO Anweisung: ein Satz über der ersten Teilaufgabe fehlt in der Bank – Nr. 17
\begin{aufgabe}{Waagerechte Tangente und Wendepunkt einer Integralfunktion an einer doppelten Nullstelle des Integranden begründen}
\begin{teile}
\teil $J(x)$ ist das Integral von $0$ bis $x$ über $f(t) = t \cdot (t - 2)^2$. Gib zwei besondere Eigenschaften des Graphen von $J$ bei $x = 2$ an und begründe. \hfill (Abitur 2017 LK)
\end{teile}
\end{aufgabe}

%% TODO Titel: Ich-kann-Satz fehlt in der Bank (Platzhalter: Kettenname) – Nr. 18
%% TODO Anweisung: ein Satz über der ersten Teilaufgabe fehlt in der Bank – Nr. 18
\begin{aufgabe}{Integralfunktion – Fehler finden, Begründen}
\begin{teile}
\teil Ich finde den Fehler: $J(x)$ ist das Integral von $0$ bis $x$ über $f(t) = t - 2$. Ole schreibt: $J$ hat die Nullstelle $x = 2$.
\teil $J(x)$ ist das Integral von $4$ bis $x$ über $f(t)$. Begründe, dass $J(4)$ immer null ist.
\end{teile}
\end{aufgabe}
